\section*{2. 爆破簇的上同调 (参见 SGA 5 VII, 8)}

\subsubsection*{2.1.}
设 $X$ 为域 $k$ 上的光滑概形, $\Delta\subset X$ 为余维 $1+r$、处处光滑的子概形. 已知 (EGA IV, 19.4.9), 存在\underline{光滑}概形 $\widetilde X$, 配有射影态射
\begin{equation}\tag{2.1.1}
f:\widetilde X\longrightarrow X,
\end{equation}
使得置
\begin{equation}\tag{2.1.2}
\widetilde\Delta=f^{-1}(\Delta)
=\Delta\times_X\widetilde X,
\end{equation}
$f$ 诱导同构
\begin{equation}\tag{2.1.3}
f|_{\widetilde X-\widetilde\Delta}:
\widetilde X-\widetilde\Delta
\xrightarrow{\ \sim\ }X-\Delta,
\end{equation}
且 $\widetilde\Delta$ 与 $\check{\mathbb P}(\underline N_{\Delta/X})$ 之间存在 $\Delta$-同构, 其中后者是 $\Delta$ 在 $X$ 中的法丛 $\underline N_{\Delta/X}$ 的直线所成的射影丛. 因而有交换图
\begin{equation}\tag{2.1.4}
\begin{tikzcd}[row sep=huge,column sep=small]
\widetilde\Delta
  \arrow[rd,"f|_{\widetilde\Delta}"']
  & \simeq &
\check{\mathbb P}(\underline N_{\Delta|X})
  \arrow[ld,"\text{典范投影}"]\\
{} & \Delta & {}
\end{tikzcd}.
\end{equation}

\subsubsection*{2.1.5.}
在特殊情形
\[
X=\operatorname{Spec}(A),\qquad
\Delta=\operatorname{Spec}(A/I),\qquad
I=(g_0,\ldots,g_r),
\]
$\widetilde X$ 可由 Hopf 构造得到 (参见 EGA IV, 19.4.11). 由下式定义从 $X-\Delta$ 到 $\mathbb P^r$ 的态射:
\[
x\in X-\Delta\longmapsto
\bigl(g_0(x),\ldots,g_r(x)\bigr)\in\mathbb P^r
\]
并将其图像 $\Gamma$ 看作嵌入 $X\times\mathbb P^r$ 中. 随后定义 $\widetilde X$ 为 $\Gamma$ 在 $X\times\mathbb P^r$ 中的概形论闭包, 态射 $f:\widetilde X\longrightarrow X$ 为复合
\[
\widetilde X\lhook\joinrel\longrightarrow
X\times\mathbb P^r\xrightarrow{\ \operatorname{pr}_1\ }X.
\]

\subsubsection*{2.1.6.}
类似地 (仍参见 EGA IV, 19.4.11), 若 $X$ 为射影概形, 且 $\Delta$ 为 $X$ 与 $1+r$ 个超曲面的交, 这些超曲面的方程记为 $G_i=0$, $i=0,\ldots,r$, 则再次考虑态射
\[
X-\Delta\longrightarrow\mathbb P^r,\qquad
x\longmapsto\bigl(G_0(x),\ldots,G_r(x)\bigr),
\]
将其图像在 $X\times\mathbb P^r$ 中的闭包取为 $\widetilde X$. 具体地, 用 $\mathbb P^r$ 上的齐次坐标 $\lambda_0,\ldots,\lambda_r$ 表示时, $\widetilde X$ 是由方程
\begin{equation}\tag{2.1.6.1}
\lambda_iG_j(X)-\lambda_jG_i(X),
\qquad 0\leq i<j\leq r.
\end{equation}
定义的 $X\times\mathbb P^r$ 的子概形. 这些方程表明, 点 $(x,\lambda)$ 属于 $\widetilde\Delta$ 当且仅当 $x\in\Delta$; 换言之, $\check{\mathbb P}(\underline N_{\Delta/X})$ 为平凡射影丛\footnote{印本把 (2.1.6.2) 中右侧的投影标为 $\operatorname{pr}_i$. 它的靶为 $\Delta$, 而在 (2.1.6.3) 中同一投影标为 $\operatorname{pr}_1$; 此处采用改正后的标记.}
\begin{equation}\tag{2.1.6.2}
\begin{tikzcd}[column sep=large,row sep=large]
\check{\mathbb P}(\underline N_{\Delta/X})
  \arrow[rr,phantom,"\simeq"]
  \arrow[dr,"\text{典范投影}"']
  & {} &
\Delta\times\mathbb P^r
  \arrow[dl,"\operatorname{pr}_1"]\\
{} & \Delta & {}
\end{tikzcd}.
\end{equation}
有如下交换图:
\begin{equation}\tag{2.1.6.3}
\begin{tikzcd}[column sep=large,row sep=large]
\Delta\times\mathbb P^r
  \arrow[r,"\sim"]
  \arrow[dr,"\operatorname{pr}_1"]
&
\widetilde\Delta
  \arrow[r,hook]
  \arrow[d,"f|_{\widetilde\Delta}"]
&
\widetilde X
  \arrow[r,hook]
  \arrow[d,"f"]
&
X\times\mathbb P^r
  \arrow[dl,"\operatorname{pr}_1"]
\\
{} & \Delta \arrow[r,hook] & X & {}
\end{tikzcd},
\end{equation}
\begin{editorialnote}
2.1.6 中用 $(G_0:\cdots:G_r)$ 定义到 $\mathbb P^r$ 的映射时, 各 $G_i$ 须为同次数的齐次方程, 即取自同一线性系. 否则改变 $x$ 的齐次坐标代表时, 该比值一般不能定义同一个射影点. 底本未在此明说这一条件; 后文 3.1 使用两个线性型, 满足此条件.
\end{editorialnote}

而复合包含映射 $\Delta\times\mathbb P^r\hookrightarrow X\times\mathbb P^r$ 由 $\Delta\hookrightarrow X$ 经基变换得到.
